% id: 85-35
% book: 베이직쎈 확률과 통계 · 05 확률변수와 확률분포
% section: 개념 25 연속확률변수와 확률밀도함수
% passage: 다음 중 $f(x)$가 $0\le X\le 2$에서 모든 실수의 값을 갖는 연속확률변수 $X$의 확률밀도함수가 될 수 있는 것은 ‘○’를, 될 수 없는 것은 ‘×’를 ( ) 안에 써넣으시오.
% figure: tikz:fig-85-35
% answer: 
% answer_source: 
% variant_level: 0 (원본 전사)
% 이 파일은 build.mjs 가 items.json 에서 만든다. 고칠 때는 items.json 을 고친다.
\smash{\raisebox{1.5mm}{\dmkichul{베이직쎈 확률과 통계 85쪽 35번}}}
\noindent\begin{minipage}[t]{0.58\linewidth}\vspace{0pt}\raggedright \hspace*{\stretch{10}}(\hspace*{2.4em})\end{minipage}\hfill\begin{minipage}[t]{0.4\linewidth}\vspace{0pt}\centering \resizebox{\ifdim\width>\linewidth\linewidth\else\width\fi}{!}{% 개념 25 85-35(인쇄 85쪽) — 확률밀도함수가 될 수 있는지 묻는 그래프: (0,0)-(1,1)-(2,0) 을 잇는 꺾은선과 점선 보조선.
% 스캔 화질이 낮아 TikZ 로 다시 그림. 축 0.7pt · 그래프 0.6pt · 라벨은 원문에 있는 것만(O · x · y · 1 · 2 · y=f(x)).
\begin{tikzpicture}[scale=0.75, line width=0.6pt]
  \draw[->, line width=0.7pt] (-0.45,0) -- (2.95,0) node[below, font=\small] {$x$};
  \draw[->, line width=0.7pt] (0,-0.6) -- (0,1.75) node[right, font=\small] {$y$};
  \node[below left=-2pt and -2pt, font=\small] at (0,0) {$\mathrm{O}$};
  \draw[dotted, line width=0.6pt] (0,1) -- (1,1) -- (1,0);
  \draw (0,0) -- (1,1) -- (2,0);
  \node[left, font=\small] at (0,1) {$1$};
  \node[below, font=\small] at (1,0) {$1$};
  \node[below, font=\small] at (2,0) {$2$};
  \node[right, font=\small] at (1.1,1.05) {$y=f(x)$};
\end{tikzpicture}
}\end{minipage}\par
